% Options for packages loaded elsewhere
\PassOptionsToPackage{unicode}{hyperref}
\PassOptionsToPackage{hyphens}{url}
%
\documentclass[
  9pt,
]{article}
\usepackage{amsmath,amssymb}
\usepackage{iftex}
\ifPDFTeX
  \usepackage[T1]{fontenc}
  \usepackage[utf8]{inputenc}
  \usepackage{textcomp} % provide euro and other symbols
\else % if luatex or xetex
  \usepackage{unicode-math} % this also loads fontspec
  \defaultfontfeatures{Scale=MatchLowercase}
  \defaultfontfeatures[\rmfamily]{Ligatures=TeX,Scale=1}
\fi
\usepackage{lmodern}
\ifPDFTeX\else
  % xetex/luatex font selection
  \setmainfont[]{DejaVu Sans}
\fi
% Use upquote if available, for straight quotes in verbatim environments
\IfFileExists{upquote.sty}{\usepackage{upquote}}{}
\IfFileExists{microtype.sty}{% use microtype if available
  \usepackage[]{microtype}
  \UseMicrotypeSet[protrusion]{basicmath} % disable protrusion for tt fonts
}{}
\makeatletter
\@ifundefined{KOMAClassName}{% if non-KOMA class
  \IfFileExists{parskip.sty}{%
    \usepackage{parskip}
  }{% else
    \setlength{\parindent}{0pt}
    \setlength{\parskip}{6pt plus 2pt minus 1pt}}
}{% if KOMA class
  \KOMAoptions{parskip=half}}
\makeatother
\usepackage{xcolor}
\usepackage{color}
\usepackage{fancyvrb}
\newcommand{\VerbBar}{|}
\newcommand{\VERB}{\Verb[commandchars=\\\{\}]}
\DefineVerbatimEnvironment{Highlighting}{Verbatim}{commandchars=\\\{\}}
% Add ',fontsize=\small' for more characters per line
\newenvironment{Shaded}{}{}
\newcommand{\AlertTok}[1]{\textcolor[rgb]{1.00,0.00,0.00}{\textbf{#1}}}
\newcommand{\AnnotationTok}[1]{\textcolor[rgb]{0.38,0.63,0.69}{\textbf{\textit{#1}}}}
\newcommand{\AttributeTok}[1]{\textcolor[rgb]{0.49,0.56,0.16}{#1}}
\newcommand{\BaseNTok}[1]{\textcolor[rgb]{0.25,0.63,0.44}{#1}}
\newcommand{\BuiltInTok}[1]{\textcolor[rgb]{0.00,0.50,0.00}{#1}}
\newcommand{\CharTok}[1]{\textcolor[rgb]{0.25,0.44,0.63}{#1}}
\newcommand{\CommentTok}[1]{\textcolor[rgb]{0.38,0.63,0.69}{\textit{#1}}}
\newcommand{\CommentVarTok}[1]{\textcolor[rgb]{0.38,0.63,0.69}{\textbf{\textit{#1}}}}
\newcommand{\ConstantTok}[1]{\textcolor[rgb]{0.53,0.00,0.00}{#1}}
\newcommand{\ControlFlowTok}[1]{\textcolor[rgb]{0.00,0.44,0.13}{\textbf{#1}}}
\newcommand{\DataTypeTok}[1]{\textcolor[rgb]{0.56,0.13,0.00}{#1}}
\newcommand{\DecValTok}[1]{\textcolor[rgb]{0.25,0.63,0.44}{#1}}
\newcommand{\DocumentationTok}[1]{\textcolor[rgb]{0.73,0.13,0.13}{\textit{#1}}}
\newcommand{\ErrorTok}[1]{\textcolor[rgb]{1.00,0.00,0.00}{\textbf{#1}}}
\newcommand{\ExtensionTok}[1]{#1}
\newcommand{\FloatTok}[1]{\textcolor[rgb]{0.25,0.63,0.44}{#1}}
\newcommand{\FunctionTok}[1]{\textcolor[rgb]{0.02,0.16,0.49}{#1}}
\newcommand{\ImportTok}[1]{\textcolor[rgb]{0.00,0.50,0.00}{\textbf{#1}}}
\newcommand{\InformationTok}[1]{\textcolor[rgb]{0.38,0.63,0.69}{\textbf{\textit{#1}}}}
\newcommand{\KeywordTok}[1]{\textcolor[rgb]{0.00,0.44,0.13}{\textbf{#1}}}
\newcommand{\NormalTok}[1]{#1}
\newcommand{\OperatorTok}[1]{\textcolor[rgb]{0.40,0.40,0.40}{#1}}
\newcommand{\OtherTok}[1]{\textcolor[rgb]{0.00,0.44,0.13}{#1}}
\newcommand{\PreprocessorTok}[1]{\textcolor[rgb]{0.74,0.48,0.00}{#1}}
\newcommand{\RegionMarkerTok}[1]{#1}
\newcommand{\SpecialCharTok}[1]{\textcolor[rgb]{0.25,0.44,0.63}{#1}}
\newcommand{\SpecialStringTok}[1]{\textcolor[rgb]{0.73,0.40,0.53}{#1}}
\newcommand{\StringTok}[1]{\textcolor[rgb]{0.25,0.44,0.63}{#1}}
\newcommand{\VariableTok}[1]{\textcolor[rgb]{0.10,0.09,0.49}{#1}}
\newcommand{\VerbatimStringTok}[1]{\textcolor[rgb]{0.25,0.44,0.63}{#1}}
\newcommand{\WarningTok}[1]{\textcolor[rgb]{0.38,0.63,0.69}{\textbf{\textit{#1}}}}
\usepackage{longtable,booktabs,array}
\usepackage{calc} % for calculating minipage widths
% Correct order of tables after \paragraph or \subparagraph
\usepackage{etoolbox}
\makeatletter
\patchcmd\longtable{\par}{\if@noskipsec\mbox{}\fi\par}{}{}
\makeatother
% Allow footnotes in longtable head/foot
\IfFileExists{footnotehyper.sty}{\usepackage{footnotehyper}}{\usepackage{footnote}}
\makesavenoteenv{longtable}
\usepackage{graphicx}
\makeatletter
\def\maxwidth{\ifdim\Gin@nat@width>\linewidth\linewidth\else\Gin@nat@width\fi}
\def\maxheight{\ifdim\Gin@nat@height>\textheight\textheight\else\Gin@nat@height\fi}
\makeatother
% Scale images if necessary, so that they will not overflow the page
% margins by default, and it is still possible to overwrite the defaults
% using explicit options in \includegraphics[width, height, ...]{}
\setkeys{Gin}{width=\maxwidth,height=\maxheight,keepaspectratio}
% Set default figure placement to htbp
\makeatletter
\def\fps@figure{htbp}
\makeatother
\setlength{\emergencystretch}{3em} % prevent overfull lines
\providecommand{\tightlist}{%
  \setlength{\itemsep}{0pt}\setlength{\parskip}{0pt}}
\setcounter{secnumdepth}{-\maxdimen} % remove section numbering
\usepackage[a4paper,margin=20mm]{geometry}
\usepackage{longtable,booktabs,array}
\usepackage{float}
\usepackage{graphicx}
\setkeys{Gin}{width=0.88\linewidth,height=0.83\textheight,keepaspectratio}
\usepackage{fancyhdr}
\pagestyle{fancy}\fancyhf{}\fancyfoot[C]{\thepage}\renewcommand{\headrulewidth}{0pt}
\setlength{\emergencystretch}{3em}
\usepackage{fvextra}
\DefineVerbatimEnvironment{Highlighting}{Verbatim}{breaklines,commandchars=\\\{\}}
\ifLuaTeX
  \usepackage{selnolig}  % disable illegal ligatures
\fi
\usepackage{bookmark}
\IfFileExists{xurl.sty}{\usepackage{xurl}}{} % add URL line breaks if available
\urlstyle{same}
\hypersetup{
  hidelinks,
  pdfcreator={LaTeX via pandoc}}

\author{}
\date{}

\begin{document}

\section{TEST PLAN --- E-COMMERCE
STORE}\label{test-plan-e-commerce-store}

\textbf{COMP.SE.200-2026-2027-1 Software Testing \textbar{} Part 1 ---
Test Plan} \textbf{Group 31 \textbar{} Pham Thanh(154888686) \textbar{}
Georgios Malezoglou (154731995)} \textbf{Version:} 1.0 \textbar{}
\textbf{Date:} 7 November 2026

\subsection{Definitions and
abbreviations}\label{definitions-and-abbreviations}

\begin{longtable}[]{@{}
  >{\raggedright\arraybackslash}p{(\columnwidth - 2\tabcolsep) * \real{0.1412}}
  >{\raggedright\arraybackslash}p{(\columnwidth - 2\tabcolsep) * \real{0.8588}}@{}}
\toprule\noalign{}
\begin{minipage}[b]{\linewidth}\raggedright
Term
\end{minipage} & \begin{minipage}[b]{\linewidth}\raggedright
Meaning
\end{minipage} \\
\midrule\noalign{}
\endhead
\bottomrule\noalign{}
\endlastfoot
SUT & System under test: the selected utility-library source files \\
E2E & End-to-end; a complete user workflow across system components \\
CI & Continuous integration; automated verification on code changes \\
REST/API & Representational State Transfer / application programming
interface \\
Coverage & Proportion of instrumented statements or branches exercised
by tests \\
Stub & Controlled substitute for a dependency during a test \\
Unit test & A test of one small unit of code in isolation \\
Quality gate & A condition that must be satisfied before testing or a
change is accepted \\
\end{longtable}

\section{1. Introduction}\label{introduction}

\subsection{1.1 Purpose and
deliverables}\label{purpose-and-deliverables}

This plan explains what will be tested and why, how the tests will be
designed, and how their results will be recorded and presented. It also
explains how the results will help on deciding whether a helper
functions library with no documentation can be reused in a food product
e-commerce shop. It includes three E2E scenarios, 10 selected library
functions in a priority order, a reproducible test approach, reporting
and defect procedures, quality gates, detailed unit-test designs, and
additional system-level testing recommendations. This is what Part 1 or
this project will deliver, while part 2 includes the implementation of
unit tests for the selected functions, GitHub Actions and Coveralls
integration {[}1{]}.

\subsection{1.2 Application context and
risk}\label{application-context-and-risk}

Customers search and filter products based on name, category, price,
ingredients and the product producer. Customers can navigate through the
results (splitted in different pages), update their cart and use an
external provider for payment. Producers can add products manually or
through an existing bulk-import portal. Inputs like name, price and
producer are mandatory while category and contents are not. The frontend
has to validate prices, descriptions and order inputs. A REST backend
takes care of third-party services and catalogue operations {[}1{]}.
Products may hide by wrong filters, totals may be affected due to
numeric conversion errors and poor validation may admit invalid product
records. The library is available in a project that has be abandoned and
documentation is missing, so a developer should not guess what each
function based on the name alone.

\section{2. Context of testing}\label{context-of-testing}

\subsection{2.1 Three representative E2E
scenarios}\label{three-representative-e2e-scenarios}

\subsubsection{S1 --- Product search, filtering and
pagination}\label{s1-product-search-filtering-and-pagination}

\textbf{Prerequisites:} Catalog has several producers, optional
attributes that could be missing, and way too many matches to be
displayed in one page. \textbf{Basic workflow:} (1) user accesses
catalog; (2) user inputs the query into the search box; (3) frontend
validates input and sends a REST request; (4) backend responds with
matches; (5) user applies filters according to category, price, content,
or producer; (6) system combines filters; (7) user navigates between
pages; (8) matches are displayed on a page based on filters applied.
\textbf{Variations:} empty/blank search, no matches, optional attributes
missing, last/incorrect page, or incorrect price interval.
\textbf{Expected output:} only valid matches are shown, paging works,
blank output is justified. \textbf{Possible utilities:} \texttt{filter},
\texttt{every}, \texttt{toNumber}, \texttt{isEmpty}, \texttt{defaultTo},
\texttt{chunk}, \texttt{slice}. \textbf{Risk:} incorrect filters or
pages numbers lead to hiding or misrepresentation of inventory. Activity
diagram: refer to Appendix A.

\subsubsection{S2 --- Cart, quantities and
checkout}\label{s2-cart-quantities-and-checkout}

\textbf{Preconditions:} The prices of products are correct; access to
the cart and the checkout. \textbf{Flow:} (1) products selected by the
customer; (2) add/remove a product to/from the cart; (3) change the
quantity; (4) calculation of subtotals and total by the cart after each
action; (5) cart overview; (6) sending checkout data from the backend to
an external provider; (7) showing the outcome on UI. \textbf{Cases to
consider:} incorrect, null, negative or non-integer quantity; empty
cart; unavailable provider; duplicate submission. \textbf{Expected
results:} accurate cart state and proper formatting of money values
without false indication of payment success. \textbf{Candidate
utilities:} \texttt{add}, \texttt{reduce}, \texttt{toNumber} (and
\texttt{toInteger}, \texttt{clamp} as alternatives). \textbf{Risk:}
incorrect totals or duplicate order. Activity Diagram: see Appendix B.
The general \texttt{add} utility is not reliable for processing decimal
money; special tests are needed for money utilities.

\subsubsection{S3 --- Producer creation and bulk
import}\label{s3-producer-creation-and-bulk-import}

\textbf{Prerequisites:} Access by the producer of the entry form or
import portal; mandatory field requirements. \textbf{Main flow:} (1)
choice of single entry or bulk import by the producer; (2) providing
product name, price, producer, and optional additional information; (3)
validation of input by the frontend/importer; (4) invalid entries
reported per the import policy; (5) valid entries stored using REST; (6)
producer verifies entries in the catalogue. \textbf{Variants:} missing
mandatory field, optional fields null values, incorrectly formatted
prices, mixed invalid/valid entries, duplicates, and bulk imports.
\textbf{Expected results:} No silent addition of invalid products,
successful import and correct catalogue. \textbf{Candidate utilities:}
\texttt{toNumber}, \texttt{isEmpty}, \texttt{every}, \texttt{defaultTo},
\texttt{countBy}; \texttt{upperFirst} is an unused candidate for
formatting purposes. \textbf{Risk:} Corrupted data in catalogue affects
search and purchase. Activity diagram: Appendix C.

\subsection{2.2 Scope and risk-based
selection}\label{scope-and-risk-based-selection}

Only ten \texttt{src} files at the highest level will be subjected to
unit test. The hidden \texttt{.internal} folder is assumed to be
error-free and hence will neither be tested nor covered, according to
the instructions {[}1{]}. The user interface for the application, REST
back-end service, and the interaction with the payment provider are
beyond the scope of the provided implementation and their suggested
tests are listed in §4.3.

The team chose the ten files above based on three criteria between 1-3
in total: \textbf{scenarios coverage} (number of S1-S3 scenarios that
can be used), \textbf{impact of failure} (consequences of a possible
error) and \textbf{estimated defect likelihood} (complexity and corner
cases). \textbf{Total = coverage + impact + likelihood.} It is not about
actual defect probability but rather an estimation for planning
purposes. In case of ties, the priority goes first by the business
impact, then by the coverage, and finally by the relevance to
implementation chosen by the group. Scores and the unselected candidates
are recorded in Appendix D.

\begin{longtable}[]{@{}
  >{\raggedright\arraybackslash}p{(\columnwidth - 4\tabcolsep) * \real{0.2133}}
  >{\raggedright\arraybackslash}p{(\columnwidth - 4\tabcolsep) * \real{0.1333}}
  >{\raggedright\arraybackslash}p{(\columnwidth - 4\tabcolsep) * \real{0.6533}}@{}}
\toprule\noalign{}
\begin{minipage}[b]{\linewidth}\raggedright
Selected file
\end{minipage} & \begin{minipage}[b]{\linewidth}\raggedright
Scenarios
\end{minipage} & \begin{minipage}[b]{\linewidth}\raggedright
Primary test focus
\end{minipage} \\
\midrule\noalign{}
\endhead
\bottomrule\noalign{}
\endlastfoot
\texttt{toNumber.js} & S1, S2, S3 & Numeric coercion, invalid values and
prices \\
\texttt{filter.js} & S1 & Predicate correctness, empty and mixed
arrays \\
\texttt{chunk.js} & S1 & Page sizes, remainder and boundary inputs \\
\texttt{isEmpty.js} & S1, S2, S3 & Empty strings, arrays, objects and
null-like data \\
\texttt{add.js} & S2 & Arithmetic, coercion and precision limitations \\
\texttt{countBy.js} & S3 & Group counts and duplicate detection
support \\
\texttt{reduce.js} & S2 & Accumulation, initial value and empty
collections \\
\texttt{every.js} & S1, S3 & All-condition predicates and empty
collections \\
\texttt{defaultTo.js} & S1, S3 & Missing optional values and valid falsy
values \\
\texttt{slice.js} & S1 & Start/end bounds and result subsets \\
\end{longtable}

The files \texttt{clamp.js}, \texttt{toInteger.js},
\texttt{toString.js}, and \texttt{upperFirst.js} were considered but not
included due to the file limit of ten. \texttt{map.js} and
\texttt{toFinite.js} are other possible choices but not included in the
final selection of the team. This is a \textbf{selection rationale} that
the team has to verify the actual exports of the functions prior to Part
2. The library helper does not perform the search, check out, or
producer validation.

\subsection{2.3 Assumptions and
limitations}\label{assumptions-and-limitations}

\texttt{.internal} functions within the library are deemed safe solely
based on the fact that the assignment allows for such an assumption to
be made, while no other \texttt{src} functions may be considered so.
Expected results need to be based on the function semantics known for
certain rather than on theoretical applications. If any of the exported
helpers forcibly manipulates the input or outputs IEEE-754 floating
point numbers, then such actions need to be mentioned along with the
potential application-based issues. There will be no application code,
backend interface, sandboxed payments, nor example products provided
{[}1{]}.

\section{3. Tools, reporting and quality
controls}\label{tools-reporting-and-quality-controls}

\subsection{3.1 Environment and
toolchain}\label{environment-and-toolchain}

The proposed \textbf{Part 2} stack is Node.js/npm, Git/GitHub,
\textbf{Jest} for test execution/assertions, \textbf{Jest's coverage
reporting} for local statement/branch data and LCOV output,
\textbf{GitHub Actions} for CI and \textbf{Coveralls} for published
coverage. This is just a planned setup. We have not yet downloaded or
used these tools. Before setting up Jest, we will look at
\texttt{package.json} to see how the library's modules are written (ESM
or CommonJS). If something does not match, it will be documented in the
final report. We will not secretly change the library's code.
\texttt{c8} is another coverage tool, but we do not need it, since Jest
already does this. Chai, Mocha and testdouble are not needed for what we
want to achieve. GitHub Actions and Coveralls will be set up in Part 2,
not now {[}1,3--6{]}.

\begin{longtable}[]{@{}
  >{\raggedright\arraybackslash}p{(\columnwidth - 4\tabcolsep) * \real{0.3333}}
  >{\raggedright\arraybackslash}p{(\columnwidth - 4\tabcolsep) * \real{0.3333}}
  >{\raggedright\arraybackslash}p{(\columnwidth - 4\tabcolsep) * \real{0.3333}}@{}}
\toprule\noalign{}
\begin{minipage}[b]{\linewidth}\raggedright
Tool
\end{minipage} & \begin{minipage}[b]{\linewidth}\raggedright
Role
\end{minipage} & \begin{minipage}[b]{\linewidth}\raggedright
Reason for choosing
\end{minipage} \\
\midrule\noalign{}
\endhead
\bottomrule\noalign{}
\endlastfoot
Node.js / npm & JavaScript runtime and dependency management. & Required
to execute the JavaScript tests and install Jest and other project
dependencies. \\
Git / GitHub & Version control, test code storage, collaboration, and
defect tracking through GitHub Issues. & Already required for the
assignment submission and provides integration with the CI pipeline. \\
Jest & Unit test execution and assertions. & The team has previous
experience with Jest. It provides a test runner, assertions, and
built-in coverage support. \\
Jest Coverage & Measure local statement and branch coverage and generate
LCOV reports. & Integrated into Jest, avoiding the need for a separate
coverage library. \\
GitHub Actions & Automatically execute tests on pushes and pull requests
(Part 2). & Integrated with GitHub and does not require a separate CI
server. \\
Coveralls & Publish coverage results and display coverage information
for individual source files (Part 2). & Accepts LCOV reports generated
by Jest and provides a centralised view of test coverage. \\
\end{longtable}

\textbf{Reproducible setup outline:} We clone the repository and write
down the commit that we tested and the Node/npm versions as well. We
look into package.json, the imports in the source files and any existing
tests. We install the dependencies (npm ci if lockfile exists, otherwise
npm install). We add Jest and write a small smoke test for one of the
selected functions and then we run it locally. We set up
collectCoverageFrom so the coverage can be measured based on the 10
selected files only and leave \textbf{/.internal/} out. We make a LCOV
report with Jest for Coveralls. In Part 2, we add a workflow to GitHub
Actions that installs the dependencies, it runs the tests and uploads
the LCOV report. We won't change library's source code and we will write
all the commands and settings to reposetory's README file. The Appendix
E shows only a small settings piece that we have not try yet.

\subsection{3.2 Execution and result
reporting}\label{execution-and-result-reporting}

For each test run record include the date/time, branch name, commit
hash, Node.js version used, test command, number of tests run, passed,
failed, skipped, blocked, planned, executed, expected vs.~actual
outputs, coverage on selected files, and links to related issues for
each test run. A test failure must be initially investigated in terms of
the test oracle and the environment before calling it a library bug.
During the test phase, keep track of the progress in a concise manner;
after the test phase, release a report on test results per
functionality, open issues, and scenarios.

\subsection{3.3 Defect tracking}\label{defect-tracking}

Use \textbf{GitHub Issues} through labels based on severity, priority,
functionality, and state. Process: \textbf{New → Confirmed →
Investigating → Resolved → Retest (if any) → Closed}; an issue other
than a defect may be \textbf{Rejected} along with the reason for that.
``Resolved'' means that the defect was resolved, although it does not
imply changing the library's code. Severity: severity/impact on the
business side (Critical -- a matter of safety in totals or lack of
function in core functionality; High -- significant incorrectness of the
result; Medium -- limited error case; Low -- minor problem with unusual
input handling). Priority: urgency (P1 -- urgent problem; P2 -- soon; P3
-- standard; P4 -- low).

\textbf{Issue Template:} Title \& Issue number; File/export that got
impacted; Test case number; Commit/Environment; Exact input and
precondition; Steps involved; Expected result (oracle); Actual result;
Logs/Screenshots (if any); Severity/Priority; Assignee \& Status; Impact
of the issue on scenario; Retest details.

\subsection{3.4 Metrics, entry/exit criteria and quality
gates}\label{metrics-entryexit-criteria-and-quality-gates}

\textbf{Entry:} ten files fixed; expected contracts reviewed; runnable
test environment and configuration setup; cases with clear inputs and
oracles. \textbf{Case pass:} actual output is equal to expected
contract, including any thrown exceptions. \textbf{Blocked:} an
unsatisfied pre-requisite stops the case execution; log separately.
\textbf{Exit:} all intended cases run or justified to be blocked;
failures analyzed and logged; coverage collected for selected files
only; uncovered branches and open risks analyzed; report generated.
Cleanliness of a stage is not reached when there are still important
failures present, despite high coverage.

\begin{longtable}[]{@{}
  >{\raggedright\arraybackslash}p{(\columnwidth - 2\tabcolsep) * \real{0.1875}}
  >{\raggedright\arraybackslash}p{(\columnwidth - 2\tabcolsep) * \real{0.8125}}@{}}
\toprule\noalign{}
\begin{minipage}[b]{\linewidth}\raggedright
Metric
\end{minipage} & \begin{minipage}[b]{\linewidth}\raggedright
Calculation / target
\end{minipage} \\
\midrule\noalign{}
\endhead
\bottomrule\noalign{}
\endlastfoot
Execution rate & Executed / planned × 100; target 100\% or documented
exceptions \\
Pass rate & Passed / executed × 100; report observed rate, do not invent
a target outcome \\
Confirmed defects & Count by function, severity and scenario
consequence \\
Statement coverage & Covered statements / measurable statements;
\textbf{aim ≥80\%} \\
Branch coverage & Covered branches / measurable branches; \textbf{aim
≥80\%} where available \\
\end{longtable}

The 80\% here refers to \textbf{group coverage} and not individual
assignment coverage. Create valid, invalid, equivalence class, boundary,
and negative test cases; examine unexecuted branches; create test cases
or justify the non-executability of some test paths. Coverage is an
indication of execution and not correctness.

\section{4. Tests and AI in testing}\label{tests-and-ai-in-testing}

\subsection{4.1 Unit-test design
principles}\label{unit-test-design-principles}

Each selected export function is tested individually using deterministic
input data and an expected outcome. Equivalence partitioning (of
valid/invalid types), boundary value analysis (with
empty/singleton/boundary page size), decision branch coverage, negative
and regression test cases must be used. Tests should be independent,
reproducible and behave-named. The validated contract rather than an
e-commerce specific feature will serve as the oracle. Implementation
detail testing and test duplication in \texttt{.internal} package should
not be included. Stubs are to be used only when there is need for
external and non-deterministic inputs; otherwise they are not required
for pure collection helper functions. In case semantics of actual source
code differs from the proposed oracle, check the contract and adjust the
test case.

\subsection{4.2 Detailed test designs for two selected
functions}\label{detailed-test-designs-for-two-selected-functions}

\subsubsection{\texorpdfstring{Function A: \texttt{chunk.js} ---
pagination support
(S1)}{Function A: chunk.js --- pagination support (S1)}}\label{function-a-chunk.js-pagination-support-s1}

\textbf{Assumed interface for test design:} \texttt{chunk(array,\ size)}
takes an array and chunks it into smaller arrays that are continuous and
consist of up to \texttt{size} elements; the last chunk can be shorter.
The actual behavior in case of a default or incorrect size value needs
to be confirmed in the repo. \textbf{No stubs} are required.

\begin{longtable}[]{@{}
  >{\raggedright\arraybackslash}p{(\columnwidth - 6\tabcolsep) * \real{0.0450}}
  >{\raggedright\arraybackslash}p{(\columnwidth - 6\tabcolsep) * \real{0.2703}}
  >{\raggedright\arraybackslash}p{(\columnwidth - 6\tabcolsep) * \real{0.4505}}
  >{\raggedright\arraybackslash}p{(\columnwidth - 6\tabcolsep) * \real{0.2342}}@{}}
\toprule\noalign{}
\begin{minipage}[b]{\linewidth}\raggedright
ID
\end{minipage} & \begin{minipage}[b]{\linewidth}\raggedright
Input
\end{minipage} & \begin{minipage}[b]{\linewidth}\raggedright
Expected output
\end{minipage} & \begin{minipage}[b]{\linewidth}\raggedright
Purpose
\end{minipage} \\
\midrule\noalign{}
\endhead
\bottomrule\noalign{}
\endlastfoot
CH-01 & \texttt{{[}1,2,3,4,5{]}}, size \texttt{2} &
\texttt{{[}{[}1,2{]},{[}3,4{]},{[}5{]}{]}} & Non-divisible page count \\
CH-02 & \texttt{{[}1,2,3,4{]}}, size \texttt{2} &
\texttt{{[}{[}1,2{]},{[}3,4{]}{]}} & Exact division \\
CH-03 & \texttt{{[}{]}}, size \texttt{3} & \texttt{{[}{]}} & Empty
catalogue \\
CH-04 & \texttt{{[}7{]}}, size \texttt{3} & \texttt{{[}{[}7{]}{]}} &
Fewer items than page size \\
CH-05 & \texttt{{[}1,2,3{]}}, size \texttt{1} &
\texttt{{[}{[}1{]},{[}2{]},{[}3{]}{]}} & Smallest positive size \\
CH-06 & \texttt{{[}1,2,3{]}}, size \texttt{5} &
\texttt{{[}{[}1,2,3{]}{]}} & Oversized page \\
CH-07 &
\texttt{{[}\textquotesingle{}a\textquotesingle{},\textquotesingle{}b\textquotesingle{},\textquotesingle{}c\textquotesingle{}{]}},
size \texttt{2} &
\texttt{{[}{[}\textquotesingle{}a\textquotesingle{},\textquotesingle{}b\textquotesingle{}{]},{[}\textquotesingle{}c\textquotesingle{}{]}{]}}
& Preserve values/order \\
CH-08 & \texttt{{[}0,null,false{]}}, size \texttt{2} &
\texttt{{[}{[}0,null{]},{[}false{]}{]}} & Preserve falsy values \\
CH-09 & \texttt{{[}1,2,3{]}}, size \texttt{0} & \textbf{Confirm source
contract before asserting} & Invalid/boundary size \\
CH-10 & \texttt{{[}1,2,3{]}}, size \texttt{-1} & \textbf{Confirm source
contract before asserting} & Negative size handling \\
\end{longtable}

CH-09/10 should be described first according to the observed contract
(without a result, coercion, error, etc.), and then the expected one
should be chosen; you should not invent a validation rule for the helper
method in your application. \textbf{Scenarios linked:} CH-01/02/04/05/06
are about first and last page rendering; S1 additionally requires tests
for saving filters between pages.

\subsubsection{\texorpdfstring{Function B: \texttt{toNumber.js} ---
numeric conversion (S1, S2,
S3)}{Function B: toNumber.js --- numeric conversion (S1, S2, S3)}}\label{function-b-tonumber.js-numeric-conversion-s1-s2-s3}

\textbf{Assumed interface for test design:} \texttt{toNumber(value)}
performs conversions according to the specified number coercion
semantics in the library; however, it \textbf{does not} check whether
the value is a correct monetary amount or has two decimal places in
currency format. Checking for \texttt{NaN} should use
\texttt{Number.isNaN}, not equality. \textbf{No stubs} will be required.
Test actual source implementation with regard to white space handling,
non-decimal formats, and objects.

\begin{longtable}[]{@{}
  >{\raggedright\arraybackslash}p{(\columnwidth - 6\tabcolsep) * \real{0.0379}}
  >{\raggedright\arraybackslash}p{(\columnwidth - 6\tabcolsep) * \real{0.1061}}
  >{\raggedright\arraybackslash}p{(\columnwidth - 6\tabcolsep) * \real{0.3788}}
  >{\raggedright\arraybackslash}p{(\columnwidth - 6\tabcolsep) * \real{0.4773}}@{}}
\toprule\noalign{}
\begin{minipage}[b]{\linewidth}\raggedright
ID
\end{minipage} & \begin{minipage}[b]{\linewidth}\raggedright
Input
\end{minipage} & \begin{minipage}[b]{\linewidth}\raggedright
Expected result
\end{minipage} & \begin{minipage}[b]{\linewidth}\raggedright
Purpose
\end{minipage} \\
\midrule\noalign{}
\endhead
\bottomrule\noalign{}
\endlastfoot
TN-01 & \texttt{\textquotesingle{}12.50\textquotesingle{}} &
\texttt{12.5} & Decimal price string \\
TN-02 & \texttt{0} & \texttt{0} & Zero remains zero \\
TN-03 & \texttt{\textquotesingle{}-3\textquotesingle{}} & \texttt{-3} &
Negative numeric conversion (not acceptance) \\
TN-04 & \texttt{\textquotesingle{}\ \ 42\ \ \textquotesingle{}} &
\texttt{42} & Surrounding whitespace \\
TN-05 & \texttt{\textquotesingle{}abc\textquotesingle{}} & \texttt{NaN}
& Malformed numeric text \\
TN-06 & \texttt{\textquotesingle{}\textquotesingle{}} & \texttt{0} &
Empty-string conversion; app must separately reject blank price \\
TN-07 & \texttt{null} & \texttt{0} & Null conversion; app must
separately reject missing price \\
TN-08 & \texttt{true} & \texttt{1} & Boolean coercion; not valid
producer input \\
TN-09 & \texttt{\textquotesingle{}0x10\textquotesingle{}} & \texttt{16}
& Nondecimal numeric notation \\
TN-10 & \texttt{\textquotesingle{}1.234\textquotesingle{}} &
\texttt{1.234} & More than two decimal places remain unrounded \\
TN-11 & \texttt{\textquotesingle{}Infinity\textquotesingle{}} &
\texttt{Infinity} & Non-finite values need separate validation \\
TN-12 & \texttt{{[}\textquotesingle{}7\textquotesingle{}{]}} &
\textbf{Confirm source contract before asserting} & Array coercion /
unexpected input type \\
\end{longtable}

\textbf{Scenario links:} TN-01/04 allow for filtering of prices and
parsing of product entries; TN-05/06/07/08/10/11 demonstrate the
necessity of validating frontend and backend after number conversion. In
S2, calculate manually the cart total (e.g., 2 × €3.25 + 1 × €1.40 =
€7.90), while checking formatting/rounding in application level,
\textbf{not} as \texttt{toNumber} promise.

\subsubsection{Coverage of the other eight selected
files}\label{coverage-of-the-other-eight-selected-files}

For \texttt{filter}, verify match/mismatch predicate, empty collections
and more than one item. For \texttt{isEmpty}, compare empty and nonempty
arrays, strings and objects. For \texttt{add}, verify positive, negative
and zero operands; floating point calculations should not use
currency-specific rounding conventions. For \texttt{countBy}, verify
repeating categories, unique entries and empty arrays. For
\texttt{reduce}, verify presence of an accumulator, empty array, and
reducer precedence. For \texttt{every}, verify matches of all, one,
none, and the empty array special case. For \texttt{defaultTo}, verify
\texttt{null} and \texttt{undefined} against \texttt{0} and
\texttt{false} (source). For \texttt{slice}, verify negative indices,
index out of range, empty result and ordering. Link test cases to
corresponding source code and S1--S3 in the Part 2 report.

\subsection{4.3 Additional testing for the complete e-commerce
system}\label{additional-testing-for-the-complete-e-commerce-system}

\textbf{Functional E2E (top business priority):} run S1 with filters, no
results, optional missing fields and pagination; S2 with quantity
changes, empty cart, totals, checkout failure and duplicate submission
protection; S3 with valid and invalid records, partial import,
duplicates and catalogue verification. Use realistic fixtures and
backend persistence verification, not only frontend notifications.

\textbf{API/integration:} validate REST API request/response schema and
errors, cart/persistences, producer API permission access, checkout
requests, retry strategy on payment provider and preventing duplicate
payments. Mock payment provider externally for deterministic failures
and run a small sandbox contract test if possible. Find out if bulk
import API call is atomic or not before validation.

\textbf{Security/data integrity:} check input validation/XSS,
authorization of producer API calls, CSRF if session based, protection
from injections, rate limiting and protection of payment data; do not
use real customer/card details. Check required fields and two decimal
places monetary policy on the server side as the frontend validation can
be easily bypassed.

\textbf{Nonfunctional and usability testing:} testing performance with
regard to big catalogues, filtering, and batch uploads; test pagination
consistency; test search/cart/errors accessibility; test compatibility
with browsers/mobility; test fault tolerance for network problems and
failure recovery. The tests related to price accuracy, authentication,
check out functionality, and data accuracy have top priority over
cosmetic tests. They are \textbf{recommendations} as the program is not
accessible {[}1{]}.

\subsection{4.4 AI use, opportunities and
safeguards}\label{ai-use-opportunities-and-safeguards}

\textbf{Actual Phase 1 use:} ChatGPT was considered in helping to draft
and edit this report, mapping scenario to function, improving the
explanation for the prioritization, and listing additional testing
risks. The responsibility for verification of all the contracts within
the library, tools used for testing, references, diagrams, and the
content of the report according to the assignment requirements still
lies with the team. Any test cases drafted with the usage of the AI are
not considered executed, no test coverage or defect results are
fabricated.

\textbf{Possible Phase 2 application:} AI (Claude) can provide
equivalence classes, edge cases, and the skeleton of tests using Jest
for the ten functions; help to analyze failing tests and uncover branch
coverage; and create well-written issue tickets. This would be
particularly useful when there is no documentation available for the
library being used and also for the edge cases of conversion/empty
inputs. Threats that need to be considered in this case would be
creating imaginary function contracts, incorrect assertion of expected
values, shallow tests which only replicate the implementation, skipping
security tests, and leaking proprietary data/code. The mitigations would
be reviewing the source code before accepting the suggestions, manually
validating the numerical oracles, keeping reviewable tests on GitHub
PRs, synthesizing the data, and making the assistance of AI transparent.

\section{5. References}\label{references}

{[}1{]} Tampere University, \emph{COMP.SE.200 Software Testing:
Practical 1 --- Part 1 Test Plan and submission instructions}, 2026
(course-provided assignment text). {[}2{]} Utility library repository:
https://github.com/otula/COMP.SE.200-2024-2025-1 {[}3{]} Jest
documentation: https://jestjs.io/docs/getting-started {[}4{]} GitHub
Actions documentation: https://docs.github.com/en/actions {[}5{]}
Coveralls documentation: https://docs.coveralls.io/ {[}6{]} Node.js
documentation: https://nodejs.org/en/learn

\newpage

\section{Appendix (excluded from ten-page main-content
limit)}\label{appendix-excluded-from-ten-page-main-content-limit}

\subsection{Appendix A --- S1 activity
diagram}\label{appendix-a-s1-activity-diagram}

\begin{figure}
\centering
\includegraphics[width=\linewidth,height=0.82\textheight,keepaspectratio]{scenario1.png}
\caption{S1 product search activity diagram}
\end{figure}

\subsection{Appendix B --- S2 activity
diagram}\label{appendix-b-s2-activity-diagram}

\begin{figure}
\centering
\includegraphics[width=\linewidth,height=0.82\textheight,keepaspectratio]{scenario2.png}
\caption{S2 cart and checkout activity diagram}
\end{figure}

\subsection{Appendix C --- S3 activity
diagram}\label{appendix-c-s3-activity-diagram}

\begin{figure}
\centering
\includegraphics[width=\linewidth,height=0.82\textheight,keepaspectratio]{scenario3.png}
\caption{S3 producer import activity diagram}
\end{figure}

\subsection{Appendix D --- Candidate prioritisation
register}\label{appendix-d-candidate-prioritisation-register}

\begin{longtable}[]{@{}
  >{\raggedright\arraybackslash}p{(\columnwidth - 12\tabcolsep) * \real{0.2537}}
  >{\raggedright\arraybackslash}p{(\columnwidth - 12\tabcolsep) * \real{0.1343}}
  >{\raggedleft\arraybackslash}p{(\columnwidth - 12\tabcolsep) * \real{0.1194}}
  >{\raggedleft\arraybackslash}p{(\columnwidth - 12\tabcolsep) * \real{0.0896}}
  >{\raggedleft\arraybackslash}p{(\columnwidth - 12\tabcolsep) * \real{0.1493}}
  >{\raggedleft\arraybackslash}p{(\columnwidth - 12\tabcolsep) * \real{0.0746}}
  >{\raggedright\arraybackslash}p{(\columnwidth - 12\tabcolsep) * \real{0.1791}}@{}}
\toprule\noalign{}
\begin{minipage}[b]{\linewidth}\raggedright
File
\end{minipage} & \begin{minipage}[b]{\linewidth}\raggedright
Scenarios
\end{minipage} & \begin{minipage}[b]{\linewidth}\raggedleft
Coverage
\end{minipage} & \begin{minipage}[b]{\linewidth}\raggedleft
Impact
\end{minipage} & \begin{minipage}[b]{\linewidth}\raggedleft
Likelihood
\end{minipage} & \begin{minipage}[b]{\linewidth}\raggedleft
Total
\end{minipage} & \begin{minipage}[b]{\linewidth}\raggedright
Decision
\end{minipage} \\
\midrule\noalign{}
\endhead
\bottomrule\noalign{}
\endlastfoot
\texttt{toNumber.js} & 1,2,3 & 3 & 3 & 2 & 8 & Selected \\
\texttt{filter.js} & 1 & 1 & 3 & 3 & 7 & Selected \\
\texttt{chunk.js} & 1 & 1 & 3 & 3 & 7 & Selected \\
\texttt{isEmpty.js} & 1,2,3 & 3 & 2 & 2 & 7 & Selected \\
\texttt{add.js} & 2 & 1 & 3 & 2 & 6 & Selected \\
\texttt{countBy.js} & 3 & 1 & 2 & 3 & 6 & Selected \\
\texttt{reduce.js} & 2 & 1 & 3 & 1 & 5 & Selected \\
\texttt{every.js} & 1,3 & 2 & 2 & 1 & 5 & Selected \\
\texttt{defaultTo.js} & 1,3 & 2 & 2 & 1 & 5 & Selected \\
\texttt{slice.js} & 1 & 1 & 2 & 2 & 5 & Selected \\
\texttt{clamp.js} & 2 & 1 & 1 & 3 & 5 & Not selected \\
\texttt{toInteger.js} & 2 & 1 & 2 & 1 & 4 & Not selected \\
\texttt{toString.js} & 1 & 1 & 1 & 2 & 4 & Not selected \\
\texttt{upperFirst.js} & 3 & 1 & 1 & 1 & 3 & Not selected \\
\end{longtable}

The score-5 tie is not uniquely resolved by impact alone: \texttt{slice}
and \texttt{clamp} both score 5, but \texttt{slice} has greater direct
relevance to the agreed pagination scenario. This documented judgement
preserves the group's ten-file decision. Scores are planning estimates,
not source-verified measurements.

\subsection{Appendix E --- Planned Jest configuration
sketch}\label{appendix-e-planned-jest-configuration-sketch}

\begin{Shaded}
\begin{Highlighting}[]
\CommentTok{// Illustrative only: confirm module format and source paths before Part 2.}
\NormalTok{module}\OperatorTok{.}\AttributeTok{exports} \OperatorTok{=}\NormalTok{ \{}
  \DataTypeTok{testEnvironment}\OperatorTok{:} \StringTok{\textquotesingle{}node\textquotesingle{}}\OperatorTok{,}
  \DataTypeTok{collectCoverage}\OperatorTok{:} \KeywordTok{true}\OperatorTok{,}
  \DataTypeTok{coverageReporters}\OperatorTok{:}\NormalTok{ [}\StringTok{\textquotesingle{}text\textquotesingle{}}\OperatorTok{,} \StringTok{\textquotesingle{}lcov\textquotesingle{}}\NormalTok{]}\OperatorTok{,}
  \DataTypeTok{collectCoverageFrom}\OperatorTok{:}\NormalTok{ [}
    \StringTok{\textquotesingle{}src/\{toNumber,filter,chunk,isEmpty,add,countBy,reduce,every,defaultTo,slice\}.js\textquotesingle{}}\OperatorTok{,}
    \StringTok{\textquotesingle{}!src/.internal/**\textquotesingle{}}
\NormalTok{  ]}
\NormalTok{\}}\OperatorTok{;}
\end{Highlighting}
\end{Shaded}

\textbf{Final submission check:} insert both student numbers; confirm
exported function contracts and adjust provisional oracles; review AI
declaration; check diagrams and page count; export the final PDF and
name it \texttt{31\_STUDENT1\_STUDENT2\_2026\_part1.pdf} with the actual
student numbers; submit to Moodle before the stated deadline {[}1{]}.

\end{document}
